\documentclass[12pt]{article}

\usepackage{lipsum}
\usepackage[margin=1in]{geometry}
\usepackage{graphicx}
\graphicspath{ {./data/} }

\usepackage{float} % for H specifier
\usepackage{subcaption} % subfigures

\usepackage{xcolor}

%\usepackage{endfloat} 
%nomarkers
%nolists
%figuresonly
%\renewcommand{\efloatseparator}{\bigskip}

%%%%%%%%%%%%%%%%%%%%%%%%%%%%%%%%%%%
\usepackage{tikz}
\newcommand{\placeholder}[2]{\begin{tikzpicture}\node[fill=red!50!yellow!30!white,minimum width=#1\linewidth,minimum height=#1\linewidth]{#2};\end{tikzpicture}}
%%%%%%%%%%%%%%%%%%%%%%%%%%%%%%%%%%%
\begin{document}

\section{Introduction}
Images in \LaTeX{} are handled by the \texttt{graphicx} package. We include the actual image file using \texttt{\textbackslash includegraphics}.

However, we usually wrap the image in a \texttt{\textbackslash begin\{figure\}} environment. This is called a \textbf{float}. \texttt{\textbackslash begin\{table\}} is another example of a float.  
Floats drift to the optimal location (top of page, bottom, etc.)to prevent ugly page breaks. You guide this behavior using \textbf{placement specifiers}:
\begin{itemize}
    \item \texttt{h} (here) - Place around here (respecting `aesthetics').
    \item \texttt{t} (top), \texttt{b} (bottom), \texttt{p} (page of floats).
    \item \texttt{H} (HERE) - Exactly here, (requires the \texttt{float} package).
\end{itemize}

\section{Single Images}
\lipsum[1-5]


\begin{figure} 
    \centering
    \placeholder{0.4}{Insert image of green cells using \texttt{includegraphics}}
    %..use \caption to add caption
    %..use \label to add a label to this figure for cross referencing
\end{figure}

{\color{red}
\lipsum[4-9]
}
%!h -> forget aesthetic put around here
% H -> put here, dont care about nothin
% h -> put around here respecting aesthetics
\begin{figure}[!t]%[!t]%[H]%[!h]%[!htbp] %[h] v/s [!h] v/s [H] push all to end coz figure voilates limits v/s
    \centering
    \includegraphics[width=\textwidth]{images/blue.jpg}
    \caption{The blue figure}
    \label{fig:blueCross}
\end{figure}

{\color{blue}
\lipsum[13-17]
}

\twocolumn
\section{Image in two column documents}
\lipsum[1]
%..insert figure in two column mode
{\color{red}\lipsum[3-9]}


\begin{figure*}[h]
    \centering
    \placeholder{0.4}{\texttt{figure*} for full figure across all columns}
    %..add caption
    %..add label
\end{figure*}
\color{blue}\lipsum[1-7]
\end{document}
